\chapter{Introducción}
\label{cap:introduccion}

\section{Motivación}

La epilepsia es una de las enfermedades neurológicas crónicas más frecuentes: afecta a unos 50 millones de personas en
el mundo \parencite{who2024}. Clínicamente se define por una predisposición duradera a generar crisis epilépticas no
provocadas; en la práctica, por dos crisis no provocadas separadas por más de 24 horas, o una crisis con un riesgo de
recurrencia alto \parencite{fisher2014}. Una \emph{crisis} es un episodio transitorio de actividad neuronal anormalmente
excesiva o sincrónica. Esa sincronía es la que hace que la crisis sea visible en el electroencefalograma (EEG): millones
de neuronas que descargan juntas producen oscilaciones de gran amplitud y frecuencia definida, muy distintas de la
actividad de fondo.

El EEG es la herramienta central para diagnosticar la epilepsia, clasificar el tipo de crisis y decidir el tratamiento.
En una unidad de monitoreo se registran pacientes durante días seguidos, con decenas de canales, para capturar sus crisis.
Revisar a mano ese volumen es lento y caro: un registro de 24 horas lleva horas de lectura de un especialista, y las
crisis ocupan una fracción mínima del tiempo (el 0,37~\% en los datos de este trabajo). Un \textbf{detector automático}
que marque los momentos sospechosos sirve para dos cosas: acortar la revisión y avisar al personal mientras la crisis
ocurre, cuando todavía se puede intervenir.

La dificultad no es separar una crisis típica de la actividad normal, que un modelo simple hace bien, sino hacerlo
\emph{con muy pocas falsas alarmas} y \emph{en un paciente que el modelo no vio}. Con 900 ventanas de 4~s por hora, una
especificidad de 99~\% da 9 falsas alarmas por hora, inaceptables en la práctica. Y el EEG de cada persona es distinto:
cambian la amplitud, las frecuencias dominantes y la forma de las crisis. Estas dos condiciones (desbalance extremo y
variabilidad entre pacientes) atraviesan todo el trabajo y son las que determinan las métricas, los modelos y las
conclusiones.

\section{El problema como aprendizaje automático}

Formalmente, el EEG de un paciente se divide en ventanas de duración fija y a cada ventana $\mathbf{X}$ se le asigna una
etiqueta $y \in \{0, 1\}$ (crisis o no). El objetivo es aprender una función $s(\mathbf{X})$, un \emph{puntaje}, tal que
la regla $\hat y = \mathds{1}[s(\mathbf{X}) > \tau]$ detecte la mayor cantidad de crisis con una tasa de falsas alarmas
acotada. Eso conecta el problema con cuatro áreas que el trabajo desarrolla:
\begin{itemize}
  \item \textbf{Procesamiento de señales.} La crisis se manifiesta como un cambio en el contenido de frecuencias. La
        transformada de Fourier y sus estimadores (periodograma, Welch, espectrograma) convierten la señal en atributos
        que hacen visible ese cambio.
  \item \textbf{Optimización.} Entrenar un modelo es minimizar una función de pérdida. Si es convexa (regresión logística,
        SVM), todo mínimo local es global y el resultado no depende del punto inicial. Si no lo es (redes neuronales), sí,
        y la elección del algoritmo y de su tasa de aprendizaje pasa a ser parte del modelo.
  \item \textbf{Optimización con restricciones.} El umbral $\tau$ no se elige minimizando el error: se elige maximizando la
        sensibilidad con una tasa de falsos positivos fija. Es un problema con restricción, que se resuelve con un
        multiplicador de Lagrange (lema de Neyman-Pearson) y, en su versión discreta, con programación lineal.
  \item \textbf{Evaluación.} Exactitud, sensibilidad, especificidad, curva ROC, AUC y matriz de confusión miden cosas
        distintas, y con clases tan desbalanceadas algunas pueden ser engañosas. A esas se suman las métricas clínicas:
        crisis detectadas, falsas alarmas por hora y latencia.
\end{itemize}

\section{Objetivos}

\begin{enumerate}
  \item Formular la detección de crisis como un problema de procesamiento de señales y de clasificación, a partir de un
        modelo matemático del fenómeno.
  \item Fundamentar matemáticamente cada paso: transformada de Fourier, convexidad y optimalidad local y global,
        multiplicadores de Lagrange y condiciones KKT, método simplex y umbral de Neyman-Pearson.
  \item Verificar numéricamente esos resultados sobre datos reales: unicidad del óptimo en un problema convexo, mínimos
        múltiples en uno no convexo, condiciones KKT en el dual del SVM.
  \item Entrenar y evaluar modelos convexos (regresión logística, SVM) y una red convolucional sobre espectrogramas.
  \item Comparar seis optimizadores de primer orden (SGD, Momentum, Nesterov, RMSprop, Adam, AdamW) sobre la misma red, con
        varias semillas y pruebas estadísticas.
  \item Evaluar con métricas clínicas: sensibilidad y especificidad por ventana, sensibilidad por evento, falsas alarmas
        por hora y latencia, sobre pacientes que no se usaron para entrenar.
\end{enumerate}

\section{Datos}

Se usan dos bases públicas, descargadas de sus fuentes oficiales:
\begin{itemize}
  \item \textbf{Bonn} \parencite{andrzejak2001}: 500 segmentos de 23,6~s de un canal, en cinco grupos (voluntarios sanos,
        pacientes fuera de crisis y pacientes durante crisis). Es chica, limpia y balanceada; sirve para estudiar el
        espectro de las crisis y para ilustrar la teoría con un problema de tres clases.
  \item \textbf{CHB-MIT} \parencite{shoeb2009,goldberger2000}: registros continuos de 23 canales de niños con epilepsia
        refractaria. Se usan 6 pacientes, 223 horas y 34 crisis. Tiene el desbalance real y permite evaluar en pacientes
        nuevos.
\end{itemize}

\section{Contribuciones}

Más allá de reproducir técnicas conocidas, el trabajo llega a algunos resultados que no eran evidentes al principio:
\begin{itemize}
  \item Una \textbf{calibración por paciente} muy simple, restar la mediana de cada atributo, se deduce del modelo
        multiplicativo del EEG y sube el AUC de la regresión logística de 0,80 a 0,93 en pacientes nuevos.
  \item En la comparativa de optimizadores, \textbf{la ventaja de SGD sobre los métodos adaptativos resultó ser un efecto de la
        tasa de aprendizaje}: con la tasa alta de cada grilla, SGD, RMSprop, Adam y AdamW quedan empatados.
  \item \textbf{Un solo paciente de validación no sirve para elegir hiperparámetros}: está saturado (AUC entre 0,996 y 0,999
        para todas las redes) y no predice el desempeño en prueba.
  \item \textbf{El modelo con menor AUC es el mejor detector}: la red convolucional tiene menor AUC que la regresión logística
        calibrada, pero detecta 11 de 14 crisis contra 7, con las mismas falsas alarmas y la mitad de latencia. La
        diferencia está en la zona de pocas falsas alarmas, que el AUC no pondera.
\end{itemize}

\section{Organización del trabajo}

\begin{itemize}
  \item El \textbf{Capítulo~\ref{cap:fenomeno}} describe el fenómeno: el origen fisiológico del EEG, un modelo matemático de
        la señal con y sin crisis, y las hipótesis que se ponen a prueba.
  \item El \textbf{Capítulo~\ref{cap:fourier}} desarrolla la transformada de Fourier discreta, el muestreo, el filtrado y la
        estimación espectral (periodograma, Welch, espectrograma), aplicados a Bonn.
  \item El \textbf{Capítulo~\ref{cap:datos}} presenta las dos bases, el preprocesamiento, el análisis exploratorio, las
        particiones y la calibración por paciente.
  \item El \textbf{Capítulo~\ref{cap:metricas}} define las métricas: matriz de confusión, sensibilidad, especificidad, curva ROC
        y AUC, métricas por evento, el umbral de Neyman-Pearson y los esquemas de validación y pruebas estadísticas.
  \item El \textbf{Capítulo~\ref{cap:optimizacion}} es el núcleo teórico: optimalidad local y global, convexidad, descenso por
        gradiente y su convergencia, multiplicadores de Lagrange y KKT, el método simplex, el dual del SVM y el lema de
        Neyman-Pearson.
  \item El \textbf{Capítulo~\ref{cap:convexos}} entrena regresión logística y SVM en Bonn y en CHB-MIT, verifica numéricamente la
        teoría y evalúa en pacientes nuevos con validación dejando un paciente afuera.
  \item El \textbf{Capítulo~\ref{cap:cnn}} describe la red convolucional, la retropropagación, los seis optimizadores con su
        análisis y el diseño del experimento.
  \item El \textbf{Capítulo~\ref{cap:resultados}} presenta la comparativa de optimizadores, la evaluación clínica de la red y su
        comparación con los modelos convexos.
  \item El \textbf{Capítulo~\ref{cap:conclusiones}} discute los resultados, las limitaciones y el trabajo futuro.
  \item El \textbf{Apéndice~\ref{ap:reproducibilidad}} explica cómo reproducir todos los resultados.
\end{itemize}

\section{Reproducibilidad}

Todo el código es propio y está en Python: cuatro notebooks, uno por etapa del trabajo, y cuatro scripts para la descarga,
el preprocesamiento, las métricas y el entrenamiento de la red. Cada tabla y cada figura del informe sale de un archivo
generado por esos notebooks. Las versiones de las bibliotecas y los pasos para reproducir están en el
Apéndice~\ref{ap:reproducibilidad}. Todo se corrió en una computadora personal sin GPU.
